% sections/migration.tex — Multi-GPU Migration
\section{Multi-GPU Migration}
\label{sec:migration}

The lifecycle of \Cref{sec:model} treats actors as bound to a single
device.  Multi-tenant deployments and rebalancing workloads need
something stronger: the ability to relocate a live actor from a
source GPU to a target GPU without losing in-flight messages,
breaking causal ordering, or violating the cross-tenant isolation of
\Cref{sec:tenancy}.  This section formalizes the migration protocol
\system{} uses for that purpose, gives a small-step semantics for
its phases, states the central safety theorem, and describes the
cross-GPU K2K extensions that the protocol relies on.  Both the
migration protocol and the cross-GPU K2K calculus are mechanized in
\TLA{}: \texttt{docs/verification/migration.tla} for the migration
state machine, \texttt{docs/verification/multi\_gpu\_k2k.tla} for the
NVLink-versus-host-stage routing.

\subsection{Phase Machine}
\label{sec:migration:phases}

A migration of actor \(a\) from source GPU \(g_s\) to target GPU
\(g_t\) proceeds through five phases:

\begin{description}
\item[\textsc{Running}] Source is authoritative; new messages land
in the source's inbox, processed normally.
\item[\textsc{Quiesced}] Source stops accepting new messages into
its inbox; new messages from senders unaware of the migration are
diverted to an in-flight buffer.
\item[\textsc{Transferring}] Source's state, inbox, and in-flight
buffer are copied to the target via a peer-to-peer transfer
(NVLink) or a host-mediated stage.
\item[\textsc{Swapping}] The routing table flips atomically from
\(g_s\) to \(g_t\); senders' next K2K \textsc{Send} resolves to
\(g_t\).
\item[\textsc{Complete}] Target is authoritative; source is inert
and may be reaped.
\end{description}

The phase variable is held in supervisor-private memory, polled by
both source and target on every tick.  The full transition
semantics, given as small-step rules, mirrors the \TLA{} actions
\texttt{BeginQuiesce}, \texttt{CaptureState}, \texttt{Transfer},
\texttt{SwapRouting}, and \texttt{Finalize} in
\texttt{migration.tla}; we give one representative rule and refer
the reader to the spec for the rest.

\begin{equation}
\inference[\textsc{SwapRouting}]
{
  \mathit{phase} = \state{Swapping} \\
  \mathit{checksum}_s = \mathit{checksum}_t
}
{
  \langle \Theta, M, t\rangle
  \trans{\mathsf{swap}(a, g_s, g_t)}
  \langle \Theta', M, t+1\rangle
}
\label{rule:migration-swap}
\end{equation}

\noindent where \(\Theta'\) updates the routing table for \(a\)
from \(g_s\) to \(g_t\) and advances \(\mathit{phase}\) to
\state{Complete}.  The premise
\(\mathit{checksum}_s = \mathit{checksum}_t\) refuses the swap if
the state captured at the source does not match what the target
restored: the state transfer is verified before the cutover, so a
corrupted transfer cannot land an inconsistent target as
authoritative.

\subsection{Safety Property}
\label{sec:migration:safety}

The safety guarantee of the migration protocol is the conjunction
of three statements: every message accepted before, during, or
after migration is eventually delivered (or in flight) at the
target; the relative order of two messages from the same source is
preserved across the cutover; and the routing table flips
atomically (no observable state in which both \(g_s\) and \(g_t\)
accept simultaneously).  We state this as a single theorem.

\begin{theorem}[Migration Safety]
\label{thm:migration-safety}
For every reachable configuration \(\langle \Theta, M, t\rangle\)
in which a migration of actor \(a\) from \(g_s\) to \(g_t\) has
been initiated:
\begin{enumerate}
\item \emph{Atomic swap}: at every step, exactly one of
\(\{g_s, g_t\}\) is the routing destination for messages addressed
to \(a\); there is no step in which both accept.

\item \emph{No message loss}: every message identifier admitted by
the system is, at every reachable configuration, in at least one
of the following sets: the source's inbox; the in-flight buffer;
the target's inbox; the delivered set.  (Source and target may
hold the same identifier between \textsc{Transfer} and
\textsc{Finalize}; this is the protocol's window in which a copy
exists on both sides for resilience against a transfer abort.
After \textsc{Finalize}, the source's representation is drained.)

\item \emph{FIFO across cutover}: for any two messages \(m_1, m_2\)
with the same source actor and \(\mathit{mid}_1 < \mathit{mid}_2\),
if both arrive in the target's inbox after \textsc{SwapRouting},
then \(m_1\) is at a lower index than \(m_2\).

\item \emph{State consistency}: at the moment of
\textsc{SwapRouting}, the target's restored state has the same
checksum as the source's captured state; the protocol refuses to
swap if they differ.
\end{enumerate}
\end{theorem}

\begin{proof}[Proof sketch]
By induction on the migration trace.  The atomic-swap clause~(1)
follows from the structure of the phase machine: the routing table
is a single variable updated only by
\textsc{SwapRouting}~(\Cref{rule:migration-swap}), and that rule
moves the value from \(g_s\) to \(g_t\) in one step;
\texttt{AcceptOnSrc} is gated by
\(\mathit{routing\_table} = \text{``src''}\) and
\texttt{AcceptOnTgt} by \(\mathit{routing\_table} = \text{``tgt''}\),
which are mutually exclusive.

The no-loss clause~(2) follows from a case analysis on which rule
applied last.  In phase \state{Running},
\texttt{AcceptOnSrc} appends to \(\mathit{src\_queue}\).  In phases
\(\{\state{Quiesced}, \state{Transferring}, \state{Swapping}\}\),
\texttt{AcceptDuringQuiesce} appends to
\(\mathit{in\_flight\_buffer}\) instead.  \texttt{Transfer}
concatenates \(\mathit{src\_queue}\) and
\(\mathit{in\_flight\_buffer}\) onto \(\mathit{tgt\_queue}\), so no
in-flight message is lost.  After \textsc{Swapping}, only
\texttt{AcceptOnTgt} can admit a message (and lands it in
\(\mathit{tgt\_queue}\)); \texttt{Finalize} drains
\(\mathit{src\_queue}\) and \(\mathit{src\_state}\) but only after
the cutover, so anything in those structures has already been
copied to the target.

The FIFO clause~(3) follows from the queue-concatenation property
of \texttt{Transfer}: \(\mathit{tgt\_queue}\) is constructed as
\(\mathit{src\_queue} \mathbin{\circ} \mathit{in\_flight\_buffer}\),
which preserves the relative order of any two messages from the
same source.  Subsequent \texttt{AcceptOnTgt} appends to the tail,
preserving order with respect to anything already in the queue.

The state-consistency clause~(4) is the precondition of
\Cref{rule:migration-swap}: \textsc{SwapRouting} is only enabled
when \(\mathit{checksum}_s = \mathit{checksum}_t\).  In the
runtime, the checksum is a CRC-32 computed over the captured state
bytes; a transfer error that flipped any bit would leave the two
checksums different, and the swap would be refused.  The
supervisor's escalation policy on a refused swap is to abort the
migration and leave the source authoritative.
\qed
\end{proof}

\Cref{thm:migration-safety} corresponds to the conjunction of the
\TLA{} invariants \texttt{AtomicSwap}, \texttt{NoMessageLoss},
\texttt{FIFOPreserved}, and \texttt{ChecksumMatch} in
\texttt{migration.tla} (lines 187--208).  The bounded model checker
under \texttt{migration.cfg} explores configurations of up to four
in-flight messages and verifies all four invariants
exhaustively.

\subsection{Cross-GPU K2K}
\label{sec:migration:cross-gpu}

The migration protocol relies on a cross-GPU extension of the K2K
calculus of \Cref{sec:k2k}.  The extension introduces two new
channel tiers:

\begin{itemize}
\item \textbf{NVLink P2P}: when source and destination GPUs are
connected by NVLink, the runtime issues a peer-to-peer copy
(\texttt{cuMemcpyPeerAsync}) directly from source-device memory to
target-device memory.  On a 2$\times$ H100~NVL system with the
12-link NVLink bundle, the runtime reaches 257.78~GB/s sustained
bandwidth at 16~MiB (81\% of the bundle's $\sim$318~GB/s
theoretical peak; \Cref{tab:mgpu-bw}) and runs 8.65$\times$ faster
than the host-staged fallback at the same size
(\Cref{tab:migration}).

\item \textbf{Host stage}: when no NVLink path exists, the runtime
falls back to a two-hop copy through host memory:
\texttt{cuMemcpyDtoHAsync} from source to a pinned host buffer,
then \texttt{cuMemcpyHtoDAsync} from the host buffer to target.
Latency is bounded by host memory bandwidth and the PCIe
round-trip; at 16~MiB the host path is bounded at $\sim$28~GB/s
effective throughput because the data crosses the PCIe + host
boundary twice.
\end{itemize}

The runtime selects the path based on a static NVLink-topology
query at startup; the choice is captured in the cross-GPU routing
table.  The relevant \TLA{} actions in
\texttt{multi\_gpu\_k2k.tla} are \texttt{RouteP2P} (NVLink path)
and \texttt{RouteHost} (host-stage fallback).

\begin{theorem}[Bounded Latency]
\label{thm:bounded-latency}
For every reachable configuration of the cross-GPU K2K calculus, no
message addressed to a destination on an NVLink-connected GPU is
ever placed on the host-stage queue.
\end{theorem}

\begin{proof}[Proof sketch]
\texttt{RouteP2P} and \texttt{RouteHost} have mutually exclusive
preconditions: \texttt{RouteP2P} requires
\texttt{HasNVLink(sg, dg)} and \texttt{RouteHost} requires
\(\neg\)\texttt{HasNVLink(sg, dg)}.  By the well-formedness of the
NVLink topology (a static input to the runtime), exactly one of the
two preconditions holds for any given pair of GPUs, so a message
takes the P2P path if and only if NVLink is available.  This is the
\texttt{BoundedLatency} invariant of
\texttt{multi\_gpu\_k2k.tla} (lines 158--168), and is preserved on
every reachable trace.
\qed
\end{proof}

\Cref{thm:bounded-latency} is operationally important: it
guarantees that adding more GPUs to a deployment never silently
degrades latency for actor pairs that should remain on the fast
path.  An NVLink topology change that severs the path between two
GPUs would cause new messages between actors on those GPUs to fall
back to the host-stage tier; the runtime exposes this transition
as a metric (\texttt{ringkernel\_k2k\_host\_stage\_messages\_total})
and the operator can observe and react.

\subsection{HLC Across Devices}
\label{sec:migration:hlc-cross}

When an actor migrates from \(g_s\) to \(g_t\), its HLC must
continue to increase monotonically (per the per-actor monotonicity
property of \Cref{prop:hlc-per-actor-monotone}).  The complication
is that \(g_s\) and \(g_t\) may have different host wall-clock
sources; even with NTP synchronization, sub-microsecond skew is
expected.  The protocol handles this by treating the source's HLC
at the moment of \textsc{CaptureState} as the retreat timestamp
for the purposes of \Cref{thm:hlc-restart-monotone}, applied at
the target.

Concretely, when the target restores the actor's state in
\textsc{Transferring}, it initializes the actor's HLC to
\(\hlc.\mathsf{merge}(h_s, h_t, w_t)\) where \(h_s\) is the
source's HLC at capture, \(h_t\) is the target's clock prior to the
restore, and \(w_t\) is the target's wall clock.  This ensures the
restored actor's HLC dominates both \(h_s\) and \(h_t\) and is
consistent with the target's wall clock.  The
\(\hlc.\mathsf{merge}\) operation is the same one used for
intra-device message receipt (\Cref{sec:supervision:hlc}); the only
unusual move is using it on a state restore rather than a message
receive.

The HLC-restart-monotonicity argument
(\Cref{thm:hlc-restart-monotone}) extends to migration with this
additional move: the proof goes through unchanged, with
\(h_{\mathrm{ret}}\) replaced by the merged value
\(\hlc.\mathsf{merge}(h_s, h_t, w_t)\), which dominates
\(h_{\mathrm{fault}} = h_s\).
